\chapter{Background}
\label{chap:background}
 
Section~\ref{sec:bg-docai} introduces document intelligence and the
structured extraction task.
Section~\ref{sec:bg-templates} examines the template-based extraction
paradigm and the failure mode that motivates this project's
template-free design requirement.
Section~\ref{sec:bg-utilitybills} establishes utility bills as the
target document class.
 
% -------------------------------------------------------------
\section{Document Intelligence and Structured Data Extraction}
\label{sec:bg-docai}
% -------------------------------------------------------------
 
Document AI --- or Document Intelligence --- covers the automated
reading, understanding, and analysis of business
documents~\cite{Cui2021DocumentAI}, drawing on natural language
processing and computer vision for tasks including layout analysis,
information extraction, document classification, and visual question
answering~\cite{Cui2021DocumentAI}.
Unstructured document data affects an estimated 95\% of organisations
at significant cost~\cite{Baviskar2021Automated}.
 
Documents differ in how machine-readable their structure is.
\textit{Structured} documents conform to a fixed schema --- database
exports, EDI formats --- and can be parsed deterministically.
\textit{Unstructured} documents contain free text with no consistent
positional organisation.
\textit{Semi-structured} documents fall between: they follow
human-readable layout conventions --- field labels, tables, groupings
--- but carry no machine-readable schema~\cite{Hamdi2021InvoiceIE}.
Processing them requires reasoning about both content and
layout~\cite{Nandi2024VDU}.
 
This project studies \textit{structured data extraction}: given a
semi-structured document, return typed values for a predefined set of
named fields.
Documents encode information through spatial layout and visual
organisation rather than explicit semantics --- they were built for
human readers, not parsers~\cite{Holt2018Invoices}.
The task therefore requires jointly interpreting text, visual features,
and spatial relationships between document
regions~\cite{Nandi2024VDU}, and automating it across diverse providers
demands techniques for handling layout variability and differing field
representations~\cite{Saout2024Invoice}.
 
% -------------------------------------------------------------
\section{The Template-Based Extraction Paradigm and Its Limits}
\label{sec:bg-templates}
% -------------------------------------------------------------
 
Structured extraction from business documents has historically been
handled by template-based systems; survey work confirms that existing
extraction techniques remain predominantly template-based or
rule-based~\cite{Baviskar2021Automated}.
The central design requirement of this project ---
\textit{template-free} generalisation across providers and languages
--- is a direct response to the failure mode described here.
 
\subsection{How Template Systems Work}
\label{subsec:bg-template-arch}
 
A template encodes per-provider layout knowledge as coordinate anchors,
keyword proximity rules, and regular expressions.
Palm~et~al.\ describe the standard architecture of commercial systems
such as Intellix~\cite{Schuster2013Intellix} and ITESOFT: fields are
manually annotated per template; incoming documents are matched to the
closest known template; candidates are extracted and scored by heuristics
such as spatial proximity and keyword
uniqueness~\cite{Palm2017CloudScan}.
Rusi\~{n}ol~et~al.\ reduce annotation cost to a single labelled example
per provider, learning a structural model from pairwise spatial
relationships between each target field and surrounding
words~\cite{Rusinol2013Templates} --- but the positional assumption
remains, and providers whose field positions vary produce significant
accuracy failures~\cite{Rusinol2013Templates}.
 
\subsection{The Generalisation Problem}
\label{subsec:bg-template-limits}
 
On familiar layouts, template systems perform well.
The problem is any new layout.
Palm~et~al.\ quantify this on 326,471 invoices: F1\,=\,0.887 on seen
layouts, F1\,=\,0.788 on
unseen~\cite{Palm2017CloudScan}.
Template-driven systems cannot reliably or automatically generalise to
unseen field layouts~\cite{Holt2018Invoices} --- and at BPAY's scale
of 26,000 registered billers, maintaining templates manually is not
feasible~\cite{Holt2018Invoices}.
Enterprise supplier populations routinely require hundreds to thousands
of templates~\cite{Krieger2021GNN}.
 
Attempts to reduce template dependence without abandoning layout-based
reasoning narrow the gap but do not close it.
Sequence-labelling systems such as CloudScan train a single model
across all layouts; however, as Krieger~et~al.\ observe, serialising
document text into word sequences discards layout information and
reintroduces positional assumptions through word
ordering~\cite{Krieger2021GNN}.
The underlying problem is that all of these approaches locate fields
by \textit{where they appear} rather than by reasoning about
\textit{what they mean}~\cite{Holt2018Invoices}.
Template-free extraction that generalises across unseen providers and
languages requires a model that reasons about field semantics from
content alone --- a capability not available in rule-based or
classical machine learning architectures.
Sections~\ref{sec:bg-llm} and~\ref{sec:bg-vlm} introduce the class of
models that provide it.
 
% -------------------------------------------------------------
\section{Utility Bills as a Document Extraction Testbed}
\label{sec:bg-utilitybills}
% -------------------------------------------------------------
 
Utility bills (electricity, gas, water) are the document class studied
in this project.
They concentrate the specific challenges that template-based systems
cannot handle: high layout variability, multilingual content, and a
mix of born-digital and scanned source material.
 
Bills and invoices are standard benchmark types in the document
intelligence literature, which is itself an indication that the
extraction problem is genuinely hard on this document class.
The SROIE dataset, introduced for the ICDAR 2019 competition on receipt
information extraction, reflects the real-world difficulties: poor
image quality, scanning distortion, folded documents, and complex
layouts~\cite{Huang2019SROIE}.
CORD covers Indonesian restaurant receipts annotated across 54 semantic
fields~\cite{Park2019CORD}.
FUNSD, comprising 199 noisy scanned forms, states that its documents
are \textit{``noisy and vary widely in
appearance''}~\cite{Jaume2019FUNSD}.
The RVL-CDIP collection of 400,000 document images across 16 categories
includes invoices as a named class~\cite{Harley2015RVLCDIP}.
 
Three properties of utility bills are directly relevant to this study.

\textit{Layout variability.}
Field positions, label conventions, and table structures differ across
providers.
Holt and Chisholm note that Australia's BPAY network alone covers
26,000 registered billers~\cite{Holt2018Invoices} ---
already enough to make per-provider templates unworkable.

\textit{Multilingual content.}
A billing period end date appears as \textit{Abrechnungszeitraum} in
German, \textit{p\'{e}riode de facturation} in French, and
\textit{periodo di fornitura} in Italian, each with different date
formats and decimal conventions.
Most document understanding benchmarks are English-only; XFUND
demonstrates this directly, noting that multilingual generalisation
remains underexplored~\cite{Xu2022XFUND}.

\textit{Document origin.}
Born-digital PDFs have accessible text layers; scanned documents
require OCR.
Kim~et~al.\ identify OCR error propagation --- recognition errors
degrading all downstream processing --- as one of three fundamental
problems with OCR-dependent
pipelines~\cite{Kim2022Donut}.

Hamdi~et~al.\ confirm that invoice extraction systems trained on known
layouts fail when new provider formats are encountered, because
adapting requires new manual
annotation~\cite{Hamdi2021InvoiceIE}.
Utility bills place all three difficulties in a single, naturally
occurring document class, making them a well-motivated testbed for
template-free extraction.